%% start of file `template.tex'.
%% Copyright 2006-2013 Xavier Danaux (xdanaux@gmail.com).
%
% This work may be distributed and/or modified under the
% conditions of the LaTeX Project Public License version 1.3c,
% available at http://www.latex-project.org/lppl/.

% Single-column layout, no icons: keeps the PDF text in reading order for ATS parsers.
\documentclass[10pt,a4paper,sans]{moderncv}
\moderncvstyle{banking}
\moderncvcolor{green}% accent color used everywhere (name, headings, rules, bullets)
\nopagenumbers{}
\usepackage[utf8]{inputenc}
\usepackage[top=0.8cm,bottom=0.8cm,left=1.5cm,right=1.5cm,footskip=0pt]{geometry}
\usepackage{fontawesome5}
% Wrap icon glyphs in an empty ActualText span so text extractors (ATS) skip them
\newcommand{\icon}[1]{\pdfliteral page{/Span <</ActualText <FEFF>>> BDC}#1\pdfliteral page{EMC}}
% Same for bullet markers, which otherwise extract on lines of their own
\let\origlabelitemi\labelitemi
\renewcommand*{\labelitemi}{\icon{\origlabelitemi}}
% Visible link: accent color plus underline, since moderncv hides link styling
\newcommand{\weblink}[2]{\href{#1}{\color{color1}\underline{#2}}}

% personal data
\name{Axel}{Mendoza}
\title{Senior MLOps Engineer}

\begin{document}

% Hand-written header instead of \makecvtitle: tighter spacing, and contact details wrap onto two lines.
% URLs are written out in full: ATS ignore link targets and only read the visible text.
\begin{center}
    {\fontsize{22}{26}\bfseries\color{color1}Axel Mendoza \color{color0}--- Senior MLOps Engineer}\\[4pt]
    {\color{color2}
    \icon{\faMapMarker*}~Mexico (Remote)\quad
    \icon{\faAt}~\href{mailto:axelmendoza@hotmail.fr}{axelmendoza@hotmail.fr}\quad
    \icon{\faLinkedin}~\href{https://www.linkedin.com/in/axelmdz}{linkedin.com/in/axelmdz}\\
    \icon{\faBlog}~\href{https://www.axelmendoza.com}{axelmendoza.com}\quad
    \icon{\faGithub}~\href{https://www.github.com/consciousml}{github.com/consciousml}}
\end{center}

\section{Experience}
\cventry{\textbf{\upshape 1 year 8 months}}{Platform Engineer, Open Source}{EKS Forge}{\mdseries Mar 2025 - Present}{}{
    \begin{itemize}
    \item Architecting and building \textbf{\weblink{https://eks-forge.readthedocs.io/latest/}{EKS Forge}}, an \textbf{open-source platform} for Kubernetes on AWS designed to be forked, understood, and operated by a small team, \textbf{deployable in about a day}
    \item \textbf{Multi-environment IaC} (dev, staging, prod) with Terraform and Terragrunt
    \item \textbf{GitOps} with ArgoCD (App of Apps) and Karpenter \textbf{node autoscaling}
    \item Private and public \textbf{ingress} with ALBs via Gateway API and the AWS Load Balancer Controller, \textbf{TLS} termination, and \textbf{DNS} via ExternalDNS
    \item \textbf{Secured network access} with a Tailscale VPN into private VPC subnets, NAT and VPC Endpoint egress, security groups, and Cilium network policies
    \item \textbf{Monitoring}, \textbf{logging}, and \textbf{alerting} with Prometheus, Grafana, and Loki
    \item \textbf{CI/CD} with GitHub Actions: Terratest suites that \textbf{deploy} and \textbf{test real clusters}, with \textbf{gated promotion} to production
    \item Published \textbf{documentation site} covering tutorials, how-tos, and architecture
    \end{itemize}
}
\cventry{\textbf{\upshape 6 years}}{Technical Writer}{ConsciousML}{\mdseries Oct 2020 - Present}{Freelance}{
    \begin{itemize}
    \item \textbf{Writing long-form articles} to bridge the gap between \textbf{Platform Engineering and ML}
    \item Worked with Neptune.ai (acquired by OpenAI), Gruntwork, the MLOps Community and my own \textbf{\weblink{https://axelmendoza.com}{blog}}
    \item Founding Terragrunt Ambassador
    \end{itemize}
}
\cventry{\textbf{\upshape 1 year 10 months}}{Senior MLOps Engineer}{Trinov}{\mdseries Apr 2023 - Jan 2025}{Paris, France}{
    \begin{itemize}
    \item Architected and built a \textbf{recommender system} that reduced operators’ time spent manually matching documents by \textbf{3$\times$}
    \item \textbf{Led a team} of 5 ML engineers
    \item Served \textbf{Transformer embeddings in real time} with BentoML and Docker on Cloud Run, backed by Vertex AI Feature Store, with CI/CD for model deployment
    \item Designed and owned \textbf{multi-environment IaC} on GCP with Terraform and Terragrunt
    \item Built \textbf{data pipelines} with Airbyte, dbt and Dagster
    \item Set up \textbf{experiment tracking} with MLflow
    \item \textbf{Monitored ML models} with OpenTelemetry and Elastic for observability
    \item \textbf{Load-tested} and \textbf{optimized} model server \textbf{inference} latency with Locust
    \end{itemize}
}
\cventry{\textbf{\upshape 2 years 3 months}}{Computer Vision Engineer}{Boxy}{\mdseries Dec 2020 - Apr 2023}{Paris, France}{
    \begin{itemize}
    \item Enabled checkout-free shopping by designing the \textbf{person tracking} and \textbf{re-identification} system with PyTorch that attributes every product to the right customer, then migrated it to \textbf{C++} to run in \textbf{real time}
    \item Reached \textbf{90\% correct product attribution} using \textbf{pose estimation} models
    \item Created a semi-automatic \textbf{annotation pipeline} to generate data for Deep Learning
    \item \textbf{Managed} and trained a team of 4 \textbf{annotators} on bounding box and pose video annotation
    \end{itemize}
}
\cventry{\textbf{\upshape 1 year 2 months}}{Computer Vision Engineer}{Siemens}{\mdseries Apr 2019 - Jun 2020}{Princeton, New Jersey, US}{
    \begin{itemize}
    \item \textbf{Co-authored a paper} at the European Heart Journal titled ``Deep learning for vessel-specific coronary artery calcium scoring''
    \item Built a coronary calcium detector with PyTorch that classifies calcium in high and low-risk arteries to support \textbf{heart disease diagnosis}
    \item \textbf{Automated} the \textbf{training} of these algorithms using Airflow
    \end{itemize}
}
\cventry{\textbf{\upshape 7 months}}{Computer Vision Intern}{Engie}{\mdseries May - Nov 2018}{Paris, France}{
    \begin{itemize}
    \item Improved security of power-plants by designing a \textbf{multi-camera vehicle re-identification and tracking} system using Keras and TensorFlow
    \end{itemize}
}
\cventry{\textbf{\upshape 5 months}}{Software Engineering Intern}{SAP}{\mdseries Feb - Jul 2016}{Paris, France}{
    \begin{itemize}
    \item Improved the quality of an Excel plugin by designing an \textbf{automatic testing platform} using SQL and Python.
    \end{itemize}
}

\section{Technical Skills}
\cvitem{Cloud \& Infra}{AWS (EKS, VPC, IAM, ALB), GCP (Vertex AI, Cloud Run, BigQuery, Pub/Sub), Terraform, Terragrunt, Docker, GitHub Actions, Linux}
\cvitem{Kubernetes}{ArgoCD, Helm, Karpenter, Cilium, Gateway API, External Secrets, ExternalDNS}
\cvitem{Observability}{Prometheus, Grafana, Loki, OpenTelemetry, Elastic}
\cvitem{ML \& Data}{PyTorch, BentoML, MLflow, Dagster, Airflow, dbt, Airbyte}
\cvitem{Languages}{Python, SQL, C++, Bash}

\section{Education}
\cventry{\textbf{\upshape 5 years}}{Software Engineering Master Degree, Data Science Major}{EPITA}{\mdseries 2013 - 2018}{Paris, France}{}

\end{document}

%% end of file `template.tex'.
